\section{Methodology: \texttt{LLMOrch-VAPT} Design}
\label{sec:methodology}

The core contribution of \texttt{LLMOrch-VAPT} is its resilient, provider-agnostic orchestration layer. This section details the system architecture, the autonomous failover mechanism, and the optimization methodologies employed to manage cost and reasoning quality.

\subsection{Master-Worker Multi-Agent Architecture}
As illustrated in Figure~\ref{fig:architecture}, \texttt{LLMOrch-VAPT} utilizes a hierarchical "Master-Worker" design. The \textbf{Master Supervisor}, implemented using LangGraph's stateful orchestration, maintains the global engagement context and manages task delegation. The \textbf{Orchestration Service} serves as the primary state machine, coordinating seven specialized \textbf{Worker Agents}.

\begin{figure}[t]
    \centering
    \resizebox{\columnwidth}{!}{
        \begin{tikzpicture}[
    node distance=1.5cm and 2.0cm,
    font=\sffamily\small,
    box/.style={draw, fill=blue!5, rounded corners, text width=2.5cm, align=center, minimum height=1cm, thick},
    agent/.style={draw, fill=green!5, rounded corners, text width=2.0cm, align=center, minimum height=0.8cm},
    provider/.style={draw, fill=orange!5, dashed, text width=1.8cm, align=center, minimum height=0.8cm},
    arrow/.style={-stealth, thick}
]

% Nodes
\node[box, fill=blue!15] (supervisor) {Master Supervisor\\(LangGraph)};
\node[box, below=of supervisor] (orchestrator) {Orchestration Service\\(Stateful Flow)};

\node[agent, left=of orchestrator] (network) {Network Agent};
\node[agent, below left=1.0cm and 0.2cm of orchestrator] (web) {Web Agent};
\node[agent, below=of orchestrator] (intel) {Vulnerability Intel};
\node[agent, below right=1.0cm and 0.2cm of orchestrator] (auth) {Auth Agent};
\node[agent, right=of orchestrator] (config) {Config Agent};

\node[box, fill=gray!10, below=2.5cm of orchestrator, text width=6cm] (protocol) {Unified Provider Protocol\\(LLMProvider Abstraction)};

\node[provider, below left=1.2cm and -0.5cm of protocol] (gemini) {Google Gemini};
\node[provider, below=1.2cm of protocol] (openai) {OpenAI GPT-4};
\node[provider, below right=1.2cm and -0.5cm of protocol] (local) {Local Ollama};

% Connections
\draw[arrow] (supervisor) -- (orchestrator);
\draw[arrow] (orchestrator) -- (network);
\draw[arrow] (orchestrator) -- (web);
\draw[arrow] (orchestrator) -- (intel);
\draw[arrow] (orchestrator) -- (auth);
\draw[arrow] (orchestrator) -- (config);

\draw[arrow] (network) .. controls +(0,-1) and +(-2,1) .. (protocol);
\draw[arrow] (web) .. controls +(0,-0.5) and +(-1,0.5) .. (protocol);
\draw[arrow] (intel) -- (protocol);
\draw[arrow] (auth) .. controls +(0,-0.5) and +(1,0.5) .. (protocol);
\draw[arrow] (config) .. controls +(0,-1) and +(2,1) .. (protocol);

\draw[arrow] (protocol) -- (gemini);
\draw[arrow] (protocol) -- (openai);
\draw[arrow] (protocol) -- (local);

% Background grouping
\begin{scope}[on background layer]
    \node[draw, gray!30, dashed, fill=gray!2, inner sep=15pt, rounded corners, fit=(network) (config)] (agent_pool) {};
    \node[below=0.1cm of agent_pool.north, gray] {\textbf{Specialized Worker Agents}};
\end{scope}

\end{tikzpicture}

    }
    \caption{System architecture of \texttt{LLMOrch-VAPT}, showing the Master-Worker hierarchy and the provider-agnostic abstraction layer.}
    \label{fig:architecture}
\end{figure}

Each worker agent is specialized for a specific security domain (e.g., Network, Web, Auth) and possesses a curated toolset, as summarized in Table~\ref{tab:agents}. This specialization allows the system to extract precise technical signals for each sub-task, which are critical for the complexity-aware routing described in Section~\ref{sec:dcat}.

\begin{table}[t]
    \centering
    \caption{Specialized Worker Agent Roles and Resource Allocation}
    \label{tab:agents}
    \resizebox{\columnwidth}{!}{
        \begin{tabular}{l l c l}
\toprule
\textbf{Agent Type} & \textbf{Primary Domain} & \textbf{Default Tier} & \textbf{Core Toolset} \\
\midrule
Network Agent & Port Scanning / Infra & Flash & nmap, masscan, ping \\
Web Agent & HTTP / OWASP Top 10 & Pro & requests, curl, bs4 \\
Intel Agent & CVE / Threat Intelligence & Flash & nvd-api, searchsploit \\
Auth Agent & Brute-force / Sessions & Reasoning & hydra, jwt-tool \\
Config Agent & Compliance / Hardening & Flash & checkov, lynis \\
API Agent & REST / GraphQL Security & Pro & postman-coll, swagger \\
Command Exec & Shell / Post-Exploit & Reasoning & terminal, metapsploit \\
\bottomrule
\end{tabular}

    }
    \vspace{1mm}
    \footnotesize{Note: Tiers are dynamically elevated by the Supervisor if the DCAT Complexity Signal exceeds $\tau_{high} = 0.85$.}
\end{table}

\subsection{Autonomous Provider Failover (APF)}
\label{sec:apf}
To address the "Operational Fragility" identified in Section~\ref{sec:intro}, we implement the \textbf{Autonomous Provider Failover (APF)} engine. The APF engine operates on the state machine logic shown in Figure~\ref{fig:apf_state}. Upon detecting a provider-side exception (e.g., HTTP 5xx, rate limiting, or reasoning timeout), the system triggers a mid-workflow checkpoint.

\begin{figure}[h]
    \centering
    \resizebox{\columnwidth}{!}{
        \begin{tikzpicture}[
    node distance=2cm,
    font=\sffamily\small,
    state/.style={draw, circle, fill=blue!5, minimum size=2.5cm, align=center, thick},
    arrow/.style={-stealth, thick, bend angle=20}
]

% States
\node[state] (exec) {S0\\Execution\\(Active Model)};
\node[state, right=4cm of exec, fill=red!5] (error) {S1\\Error Detection\\(Timeout/Rate/5xx)};
\node[state, below=3cm of error, fill=orange!5] (switch) {S2\\Failover Logic\\(Backend Swap)};
\node[state, below=3cm of exec, fill=green!5] (resume) {S3\\State Resume\\(Checkpointing)};

% Transitions
\draw[arrow] (exec) -- node[above] {Exception} (error);
\draw[arrow] (error) -- node[right] {Trigger APF} (switch);
\draw[arrow] (switch) -- node[above] {Retrieve State} (resume);
\draw[arrow] (resume) -- node[left] {Restart Chain} (exec);

% Self-loops
\draw[arrow, loop above] (exec) to node {Valid Response} (exec);

% Labels
\node[right=0.2cm of switch, gray, text width=3cm] {\textit{MTTR < 2.0s\\No state loss}};
\node[left=0.2cm of resume, gray, text width=3cm] {\textit{Persistence via\\PostgreSQL/Qdrant}};

\end{tikzpicture}

    }
    \caption{State machine logic of the Autonomous Provider Failover (APF) engine.}
    \label{fig:apf_state}
\end{figure}

The \texttt{LLMProvider} abstraction layer then re-instantiates the agentic chain with the next available provider in the hierarchy $\mathcal{P}$, restoring the reasoning state from the persistence layer (PostgreSQL). This mechanism ensures that the engagement resumes from the last valid reasoning step with minimal overhead, targeting a Mean Time To Recovery (MTTR) of less than 2 seconds.

\subsection{Dynamic Complexity-Aware Tiering (DCAT)}
\label{sec:dcat}
The \textbf{Dynamic Complexity-Aware Tiering (DCAT)} engine optimizes the economic sustainability of the engagement. As shown in the workflow in Figure~\ref{fig:dcat_flow}, the system analyzes each task $t$ to extract technical signals.

\begin{figure}[t]
    \centering
    \resizebox{\columnwidth}{!}{
        \begin{tikzpicture}[
    node distance=1.2cm and 1.5cm,
    font=\sffamily\small,
    startstop/.style={draw, fill=red!10, rounded corners, minimum height=0.8cm, text width=2.5cm, align=center},
    process/.style={draw, fill=blue!5, minimum height=0.8cm, text width=3cm, align=center},
    decision/.style={draw, diamond, fill=yellow!10, aspect=2, minimum height=0.8cm, text width=2.5cm, align=center, inner sep=0pt},
    arrow/.style={-stealth, thick}
]

% Nodes
\node[startstop] (input) {Input Security Task\\(User/Agent)};
\node[process, below=of input] (analysis) {Signal Extraction\\(Keywords, CVEs, Depth)};
\node[process, below=of analysis] (scoring) {Complexity Scoring\\(Eq. \ref{eq:complexity})};
\node[decision, below=of scoring] (complexity) {Complexity?};

\node[process, below left=1.5cm and 0.5cm of complexity, fill=green!10] (flash) {Flash Tier\\(e.g., Gemini Flash)};
\node[process, below=2cm of complexity, fill=blue!15] (pro) {Pro Tier\\(e.g., GPT-4o)};
\node[process, below right=1.5cm and 0.5cm of complexity, fill=purple!10] (reasoning) {Reasoning Tier\\(e.g., O1-preview)};

\node[process, below=4cm of complexity] (execution) {Task Execution\\(Managed Callback)};
\node[startstop, below=of execution] (output) {Synthesized Result};

% Paths
\draw[arrow] (input) -- (analysis);
\draw[arrow] (analysis) -- (scoring);
\draw[arrow] (scoring) -- (complexity);

\draw[arrow] (complexity) -| node[above, pos=0.5] {Low} (flash);
\draw[arrow] (complexity) -- node[right] {Moderate} (pro);
\draw[arrow] (complexity) -| node[above, pos=0.5] {High} (reasoning);

\draw[arrow] (flash) |- (execution);
\draw[arrow] (pro) -- (execution);
\draw[arrow] (reasoning) |- (execution);
\draw[arrow] (execution) -- (output);

% Labels
\node[right=0.5cm of scoring, gray, text width=3cm] {\textit{Signals: CVE-ID, Network Size, Auth Level}};

\end{tikzpicture}

    }
    \caption{Workflow of the Dynamic Complexity-Aware Tiering (DCAT) methodology.}
    \label{fig:dcat_flow}
\end{figure}

We formalize the \textbf{Complexity Signal} $C(t)$ in Equation~\ref{eq:complexity}:
\begin{equation}
\label{eq:complexity}
\begin{split}
C(t) = \frac{1}{\mathcal{K}_{max}} \left( \sum_{i \in \mathcal{K}_t} w_i \cdot \mathbb{I}(k_i \in t) \right) \\
+ \alpha \cdot \mathbb{I}(t \cap \mathcal{V}_{cve} \neq \emptyset) + \beta \cdot \mathcal{D}_{net}
\end{split}
\end{equation}
where:
\begin{itemize}
    \item $C(t) \in [0, 1]$ is the normalized complexity signal for task $t$.
    \item $w_i$ is the weight of security keyword $k_i$.
    \item $\mathbb{I}(\cdot)$ is the indicator function.
    \item $\alpha, \beta$ are scaling factors for vulnerability presence and network depth $\mathcal{D}_{net}$.
\end{itemize}


The \textbf{Optimization Manager} then solves for the optimal model $m \in \mathcal{M}$ that maximizes reasoning quality $Q$ within the constraints of the engagement budget and latency requirements, as formalized in Equation~\ref{eq:optimization}:
\begin{equation}
\label{eq:optimization}
\mathcal{J} = \arg \max_{m \in \mathcal{M}} \left\{ \mathbb{E}[Q(m, t)] \cdot \mathcal{P}(m \text{ is avail}) \right\}
\end{equation}
subject to:
\begin{equation*}
\begin{aligned}
\text{Cost}(m, t) &\leq \mathcal{B}_{task} \\
\text{Latency}(m, t) &\leq \mathcal{L}_{max} \\
\text{SecurityTier}(m) &\geq C(t)
\end{aligned}
\end{equation*}
where $Q(m, t)$ is the predicted reasoning quality of model $m$ on task $t$, and $\mathcal{B}_{task}$ is the allocated token budget.


By routing reconnaissance and routine tool-interpretation tasks to "Flash" tiers while reserving "Pro" and "Reasoning" tiers for exploitation and chain-synthesis, DCAT achieves significant cost reductions without sacrificing accuracy.

\subsection{Security-Semantic Caching (SSC)}
\label{sec:ssc}
To improve scalability across large-scale heterogeneous networks, we implement \textbf{Security-Semantic Caching (SSC)}. Unlike generic text caches, SSC utilizes the mechanism shown in Figure~\ref{fig:ssc_mechanism} to store and retrieve the \textit{reasoning patterns} behind specific vulnerability identifications.

\begin{figure}[h]
    \centering
    \resizebox{\columnwidth}{!}{
        \begin{tikzpicture}[
    node distance=1.5cm,
    font=\sffamily\small,
    box/.style={draw, fill=blue!5, rounded corners, minimum width=3cm, minimum height=1cm, align=center},
    db/.style={draw, cylinder, fill=yellow!5, shape border rotate=90, minimum width=2cm, minimum height=2.5cm, align=center},
    arrow/.style={-stealth, thick}
]

% Nodes
\node[box] (query) {Security Query\\(e.g., "Scan vuln X")};
\node[box, right=of query] (embed) {Embedding Gen\\(BAAI/bge-base)};
\node[db, right=of embed] (vector) {Vector Store\\(Qdrant)};
\node[box, below=2cm of vector] (judge) {Similarity Judger\\($\sigma \geq \tau_{sim}$)};

\node[box, left=of judge, fill=green!10] (hit) {Cache HIT\\(Return Reasoning)};
\node[box, below=1.5cm of judge, fill=red!10] (miss) {Cache MISS\\(Route to LLM)};

% Connections
\draw[arrow] (query) -- (embed);
\draw[arrow] (embed) -- (vector);
\draw[arrow] (vector) -- (judge);
\draw[arrow] (judge) -- node[above] {Yes} (hit);
\draw[arrow] (judge) -- node[right] {No} (miss);

\draw[arrow] (miss.east) .. controls +(3,0) and +(3,-2) .. (vector.south);
\node[right=2.5cm of miss, gray] {Store New Reasoning};

% Labels
\node[above=0.2cm of vector] {\textbf{SSC Memory Layer}};

\end{tikzpicture}

    }
    \caption{Architectural design of the Security-Semantic Caching (SSC) layer.}
    \label{fig:ssc_mechanism}
\end{figure}

The similarity between a new security query $q_{new}$ and cached reasoning $q_{cached}$ is determined by the cosine similarity of their vector embeddings $\mathbf{e}$, as defined in Equation~\ref{eq:similarity}:
\begin{equation}
\label{eq:similarity}
\sigma(\mathbf{e}_q, \mathbf{e}_{cache}) = \frac{\mathbf{e}_q \cdot \mathbf{e}_{cache}}{\|\mathbf{e}_q\| \|\mathbf{e}_{cache}\|} \geq \tau_{sim}
\end{equation}
where $\mathbf{e}_q$ and $\mathbf{e}_{cache}$ are vector embeddings of the current query and the cached entry, and $\tau_{sim}$ is the similarity threshold (typically 0.95 for security reasoning).
\begin{equation}
\text{Cache}(q) = 
\begin{cases} 
R_{cache} & \text{if } \sigma \geq \tau_{sim} \\
\text{LLM}(q) & \text{otherwise}
\end{cases}
\end{equation}


This allow \texttt{LLMOrch-VAPT} to bypass redundant LLM reasoning for identical vulnerability signatures across different hosts, further optimizing the operational overhead.
